This paper introduced KCH-MedRank, a knowledge-constrained learning-to-rank framework for biomedical evidence retrieval. The central finding is that medical knowledge constraints turn hypergraph structure from noise into signal. Raw hypergraph topology applied without domain-grounded concept annotations degrades ranking. Add MeSH, entity, and PrimeKG constraints and the same architecture produces interpretable, statistically significant improvements.

On BioASQ, KCH-MedRank significantly outperforms a true MedCPT Cross-Encoder baseline in Recall@10 ($+0.0157$, $p=0.0076$) and runs $18.6\times$ faster at reranking time. Flat biomedical knowledge features carry most of the aggregate gains. The hypergraph contributes targeted early-rank benefit at $k{=}5$, concentrated in questions that lack direct MeSH overlap between query and evidence. BEIR NFCorpus confirms the robustness of supervised retrieval-feature reranking, and the PubMedQA diagnostic shows that evidence ordering affects downstream answer selection under controlled conditions.

The structural ablation gives direct evidence that domain-grounded concept annotations are necessary for graph-based biomedical retrieval. Generic GraphRAG and HyperGraphRAG pipelines that build edges from lexical or embedding similarity risk introducing structural noise rather than useful signal. Interpretability analysis traces 648 rescued gold passages to three knowledge-grounded mechanisms; MeSH hierarchy paths account for 75.9\% of rescues. Error analysis reveals that biomedical concept normalization, in particular handling synonyms and abbreviations, is the most impactful direction for improving the method.

KCH-MedRank is an interpretable and efficient evidence reranking module. The code and processed feature files will be released upon publication.
